\documentclass{article}
\usepackage[utf8]{inputenc}
\usepackage[T1]{fontenc}
\usepackage{geometry}
\geometry{
    a4paper,
    left=2.5cm,
    right=2.5cm,
    top=2.5cm,
    bottom=2.5cm
}
\usepackage{graphicx}
\usepackage{float}
\usepackage{hyperref}
\usepackage[english]{babel}
\usepackage{amsmath}
\usepackage{amssymb}
\usepackage{array}
\usepackage{booktabs}
\usepackage{caption}
\usepackage{subcaption}
\usepackage{enumitem}
\usepackage{fancyhdr}
\usepackage{titlesec}
\usepackage{url}
\usepackage{textcomp}
\usepackage{xcolor}
\usepackage{listings}
\usepackage{multirow}
\usepackage{longtable}
\usepackage{tabularx}
\usepackage{times}

% Set up page style
\pagestyle{fancy}
\fancyhf{}
\rfoot{\thepage}
\renewcommand{\headrulewidth}{0pt}

% Title formatting
\titleformat{\chapter}[display]{\normalfont\bfseries}{\chaptertitle\thechapter}{0pt}{\Huge}
\titlespacing*{\chapter}{0pt}{50pt}{40pt}

% Document info
\title{Template Document}
\author{Academic Support Agent}
\date{\today}

\begin{document}
\maketitle

\maketitle

\begin{abstract}
xxx xxxxx xxxxx xxx YYYYY $\alpha_x^2+\beta_y^2=r_z^2$ xxxxx xxxxx xxxx xxxxx xxxx xxxxxxxxx
xxx xxxxx xxxxx xxx xxxx xxxx xxxxx xxxxx xxxx xxxxx xxxx xxxxxxxxx
xxx xxxxx xxxxx xxx xxxx xxxx xxxxx xxxxx xxxx xxxxx xxxx xxxxxxxxx
xxx xxxxx xxxxx xxx xxxx xxxx xxxxx xxxxx xxxx xxxxx xxxx xxxxxxxxx
\end{abstract}



\section{Introduction}
The purpose of this experiment is to determine 
xxx xxxxx xxxxx xxx xxxx xxxx xxxxx xxxxx xxxx xxxxx xxxx xxxxxxxxx
xxx xxxxx xxxxx xxx xxxx xxxx xxxxx xxxxx xxxx xxxxx xxxx xxxxxxxxx
 
xxx xxxxx xxxxx xxx xxxx xxxx xxxxx xxxxx xxxx xxxxx xxxx xxxxxxxxx
xxx xxxxx xxxxx xxx xxxx xxxx xxxxx xxxxx xxxx xxxxx xxxx xxxxxxxxx
xxx xxxxx xxxxx xxx xxxx xxxx~\cite{AG:82} xxxx xxxxx xxxx xxxxxxxxx 
xxx xxxxx xxxxx xxx xxxx xxxx xxxxx xxxxx xxxx xxxxx xxxx xxxxxxxxx
This is the way to insert a figure or simply 
leave some white space for a figure that is to be 
pasted in later, like a photo or a hand-drawn sketch.  As seen in
Figure~\ref{fig:picture}, 
\begin{figure}
  \centerline{\includegraphics{einstein.eps}}
  \caption{This is the caption for the picture.}
  \label{fig:picture}
\end{figure}
everything is clear.
xxx xxxxx xxxxx xxx xxxx xxxx xxxxx xxxxx xxxx xxxxx xxxx xxxxxxxxx
xxx xxxxx xxxxx xxx xxxx xxxx xxxxx xxxxx xxxx xxxxx xxxx xxxxxxxxx
xxx xxxxx xxxxx xxx xxxx xxxx xxxxx xxxxx xxxx xxxxx xxxx xxxxxxxxx

xxx xxxxx xxxxx xxx xxxx xxxx xxxxx xxxxx xxxx xxxxx xxxx xxxxxxxxx
xxx xxxxx xxxxx xxx xxxx xxxx xxxxx xxxxx xxxx xxxxx xxxx xxxxxxxxx
Text before the footnote.\footnote{Here's the text of the footnote.}
Text after the footnote. xxxx xxxxx xxxx xxxxx xxxxx xxxx xxxxxxxxx
xxx xxxxx xxxxx xxx xxxx xxxx xxxxx xxxxx xxxx xxxxx xxxx xxxxxxxxx

\section{Experimental procedure}

The experimental procedure described in the lab manual~\cite[pp.~13--17]{manual}
was followed exactly.  We found that the following changes were necessary to
make the circuit work:
\begin{itemize}
\item 
xxx xxxxx xxxxx xxx xxxx xxxx xxxxx xxxxx xxxx xxxxx xxxx xxxxxxxxx
xxx xxxxx xxxxx xxx xxxx xxxx xxxxx xxxxx xxxx xxxxx xxxx xxxxxxxxx.
\item
xxx xxxxx xxxxx xxx xxxx xxxx xxxxx xxxxx xxxx xxxxx xxxx xxxxxxxxx
xxx xxxxx xxxxx xxx xxxx xxxx xxxxx xxxxx xxxx xxxxx xxxx xxxxxxxxx.
The modified circuit diagram is presented in Figure~\ref{fig:circuit}.
\begin{figure}[htbp]
  \centerline{\includegraphics[width=0.8\textwidth]{circuit.eps}}
  \caption{The modified circuit diagram was drawn with {\tt edgr} and
           saved as a PostScript file.}
  \label{fig:circuit}
\end{figure}
\item
xxx xxxxx xxxxx xxx xxxx xxxx xxxxx xxxxx xxxx xxxxx xxxx xxxxxxxxx
xxx xxxxx xxxxx xxx xxxx xxxx xxxxx xxxxx xxxx xxxxx xxxx xxxxxxxxx.
\end{itemize}


\section{Results}

xxx xxxxx xxxxx xxx xxxx xxxx xxxxx xxxxx xxxx xxxxx xxxx xxxxxxxxx 
xxx xxxxx xxxxx xxx xxxx xxxx xxxxx xxxxx xxxx xxxxx xxxx xxxxxxxxx 
xxx xxxxx xxxxx xxx xxxx xxxx xxxxx xxxxx xxxx xxxxx xxxx xxxxxxxxx
xxx xxxxx xxxxx xxx xxxx xxxx as shown in Equation~\ref{eq:power}
\begin{equation}
  P = I V = I^2 R = { {V^2} \over R } \quad ,
  \label{eq:power}
\end{equation}
where $V$ is voltage~(V), $I$ is current~(A), and 
xxxxx xxxxx xxxx xxxxx xxxx xxxxxxxxx.

xxx xxxxx xxxxx xxx xxxx xxxx xxxxx xxxxx xxxx xxxxx xxxx xxxxxxxxx
xxx xxxxx xxxxx xxx xxxx xxxx xxxxx xxxxx xxxx xxxxx xxxx xxxxxxxxx
Multi--line equations are also possible (see Equation~\ref{eq:multi})
\begin{eqnarray}
  P & = & I V \nonumber \\
    & = & I^2 R
\label{eq:multi}
\end{eqnarray}
where some lines may be left unnumbered.
xxx xxxxx xxxxx xxx xxxx xxxx xxxxx xxxxx xxxx xxxxx xxxx xxxxxxxxx
xxx xxxxx xxxxx xxx xxxx xxxx xxxxx xxxxx xxxx xxxxx xxxx xxxxxxxxx.

The summary of the experimental results is presented in Table~\ref{tab:summary}
\begin{table}[htp]
\caption{This is the caption for the summary table, usually placed at the top,
         not at bottom as for figures.}
\label{tab:summary}
\begin{center}
\begin{tabular}{ccc}
\hline
 $P$, W      & $V$, V       & $I$, A 	\\
\hline
$0.0\pm0.01$ & $0.0\pm0.01$ & $0.0\pm0.01$	\\
 0.1         & 0.2          &  0.3       	\\
 0.1         & 0.2          &  0.3       	\\
 0.1         & 0.2          &  0.3       	\\
 0.1         & 0.2          &  0.3       	\\
 0.1         & 0.2          &  0.3       	\\
\hline
\end{tabular}
\end{center}
\end{table}
xxxx xxx xxxx xxxx xxxxx xxxxx xxxx xxxxx
xxxx xxx xxxx xxxx xxxxx xxxxx xxxx xxxxx
xxxx xxx xxxx xxxx xxxxx xxxxx xxxx xxxxx.

xxx xxxxx xxxxx xxx xxxx data points with error bars, and the
curve fitted to the data using Equation~\ref{eq:power} are shown
in Figure~\ref{fig:plot}
\begin{figure}[htp]
  \centerline{\includegraphics[width=0.8\textwidth]{plot.eps}}
  \caption{This plot was drawn using {\tt physica} and saved as a PostScript file.} 
  \label{fig:plot} 
\end{figure}
and the text after the figure.  The raw data is presented in Appendix~\ref{app:data}.
xxxx xxxxx xxxxx xxxx xxxxx xxxx xxxxxxxxx.

The graph of Figure~\ref{fig:plot} was generated using {\tt physica} 
(see the macro file shown in 
Appendix~\ref{app:macro}).  The data points are shown as well
as the curve representing the least--squares fit to the data.
{\tt physica} reports the value of the slope of the graph to be
$-21452\pm801$, which is equal to $h\nu/k$.  Thus from
the graph one can determine everything.

\subsection{Sample Calculation for Power}
The following calculation shows how the power was obtained.  Note that this value is
proportional to power, but it not the value for power itself.
\begin{eqnarray*}
  P & \propto & V^2 \\
    & \propto & (0.0017V-0.0004V)^2 \\
    & \propto & 1.69\times10^{-6} V^2
\end{eqnarray*}

\section{Discussion and Conclusions}
The results obtained in this experiment are quite accurate.  The graph 
presented in Figure~\ref{fig:plot}
xxx xxxxx xxxxx xxx xxxx xxxx xxxxx xxxxx xxxx xxxxx xxxx
xxxxxxxxx xxx xxxxx xxxxx xxx xxxx xxxx xxxxx xxxxx xxxx xxxxx xxxx xxxxxxxxx
xxx xxxxx xxxxx xxx xxxx xxxx xxxxx xxxxx xxxx xxxxx xxxx xxxxxxxxx
xxx xxxxx xxxxx xxx xxxx xxxx xxxxx xxxxx xxxx xxxxx xxxx xxxxxxxxx

Possible sources of error 
xxx xxxxx xxxxx xxx xxxx xxxx xxxxx xxxxx xxxx xxxxx xxxx
xxxxxxxxx xxx xxxxx xxxxx xxx xxxx xxxx xxxxx xxxxx xxxx xxxxx xxxx xxxxxxxxx
xxx xxxxx xxxxx xxx xxxx xxxx xxxxx xxxxx xxxx xxxxx xxxx xxxxxxxxx
xxx xxxxx xxxxx xxx xxxx xxxx xxxxx xxxxx xxxx xxxxx xxxx xxxxxxxxx

In conclusion, this experiment 
xxx xxxxx xxxxx xxx xxxx xxxx xxxxx xxxxx xxxx xxxxx xxxx
xxxxxxxxx xxx xxxxx xxxxx xxx xxxx xxxx xxxxx xxxxx xxxx xxxxx xxxx xxxxxxxxx
xxx xxxxx xxxxx xxx xxxx xxxx xxxxx xxxxx xxxx xxxxx xxxx xxxxxxxxx
xxx xxxxx xxxxx xxx xxxx xxxx xxxxx xxxxx xxxx xxxxx xxxx xxxxxxxxx


\begin{thebibliography}{9}

\bibitem{AG:82}
A.~Green.
\newblock {\it Reviews of Scientific Instruments}, {\bf 55}:123--134, 1982.

\bibitem{manual}
E.~Sternin.
\newblock {\it Experimental Physics II (Electronics)}, Laboratory Manual.
Brock University, 1994.

\end{thebibliography}




\section{Raw data}
\label{app:data}

\begin{table}[h]
\caption{Resistance and Temperature of the Filament}
\label{tab:data}
\vspace{0.15in}
\begin{center}
\begin{tabular}{|c|c|c|c|}
\hline
$R(T)$, $\Omega$ & $T$, K & $1/T$, K$^{-1}$ & $\ln P$ \\
\hline
151.00$\pm$3.92 & 828.35$\pm$23.46& $1.2072\times10^{-3}$& -13.29 \\
157.12$\pm$3.71 & 856.88$\pm$22.25& $1.1671\times10^{-3}$& -12.64 \\
162.53$\pm$3.49 & 881.99$\pm$21.02& $1.1338\times10^{-3}$& -12.33 \\
166.67$\pm$3.33 & 901.14$\pm$20.13& $1.1097\times10^{-3}$& -11.90 \\
171.84$\pm$3.17 & 924.98$\pm$19.25& $1.0811\times10^{-3}$& -11.25 \\ 
176.84$\pm$3.04 & 947.96$\pm$18.53& $1.0549\times10^{-3}$& -10.77 \\
181.46$\pm$2.90 & 969.13$\pm$15.49& $1.0319\times10^{-3}$& -10.20 \\
186.49$\pm$2.79 & 992.09$\pm$17.18& $1.0080\times10^{-3}$& -9.66 \\
190.91$\pm$2.69 & 1012.21$\pm$16.65& $9.8794\times10^{-4}$& -9.13 \\
195.48$\pm$2.59 & 1032.95$\pm$16.45& $9.6811\times10^{-4}$& -8.60 \\
199.93$\pm$2.50 & 1053.08$\pm$15.65& $9.4960\times10^{-4}$& -8.10 \\
204.47$\pm$2.41 & 1073.56$\pm$15.19& $9.3148\times10^{-4}$& -7.63 \\
208.62$\pm$2.34 & 1092.22$\pm$14.83& $9.1556\times10^{-4}$& -7.16 \\
\hline
\end{tabular}
\end{center}  
\end{table}

\newpage
\section{A {\tt physica} macro}
\label{app:macro}
This {\tt physica} macro was used to generate the plot of Figure~\ref{fig:plot}
as well as to fit xxx xxxxx xxxxx xxx xxxx xxxx xxxxx xxxxx xxxx xxxxx xxxx xxxxxxxxx
xxx xxxxx xxxxx xxx xxxx xxxx xxxxx xxxxx xxxx xxxxx xxxx xxxxxxxxx

{\tt 
\begin{verbatim}

      ! exp_3.pcm
      clear
      
      ! read in the data
      read\format\noerror exp_3.dat (*) x,y,dy

      ! plot the data
      label\x `Voltage, V'
      label\y `Power, W'
      set colour 1 1
      set pchar -4
      graph x,y,dy

      ! fit and plot the curve
      scalar\vary A,T,w,phi
      ! initial values for parameters
      A=2.3
      w=6.5
      phi=0
      T=10.

      fit y=A*cos(w*x+phi)*exp(-x**2/T)
      fit\update f
      set colour 2 2
      set pchar 0
      graph\noaxes x,f

\end{verbatim}
}
\end{document}